\SourceLeaf{535}\par \textsc{Gauss} collects in the following overview the methods used so far by calculators to determine $\frac{\pi}{4}=(1)$
\begingroup\fontsize{10}{13}\selectfont\setlength{\tabcolsep}{3pt}
\begin{longtable}{@{}l p{.74\linewidth}@{}}
\textsc{Machin} & $(1)=4(5)-(239)$\par also \textsc{Clausen} \\[3pt]
\textsc{Euler} & ${}=(2)+(3)$\par (\textsc{Euler} to \textsc{Goldbach} 1746 May 28) \\[3pt]
\textsc{Vega} & ${}=5(7)+2(\frac{79}{3})$\par (\textsc{Vega} Thesaurus logar. p. 633) \\[3pt]
\textsc{Vega} & ${}=2(3)+(7)$\par also \textsc{Clausen} (Astr. Nachr. vol. 25. p. 209) \\[3pt]
\textsc{Rutherford} & ${}=4(5)-(70)+(99)$\par (Philos. Trans. 1841. p. 283) \\[3pt]
\textsc{Dase} & ${}=(2)+(5)+(8)$\par (\textsc{Crelle} Journal. vol. 27. p. 198) \\[3pt]
\textsc{Gauss}. 1. & ${}=12(18)+8(57)-5(239)$ \\[3pt]
\textsc{Gauss}. 2. & ${}=12(38)+20(57)+7(239)+24(268)$ \\[3pt]
\end{longtable}\endgroup
\par The first calculations for the tables date from the preparation of the \textit{Disquiss. Arr.}; they were advanced especially in 1846 and 47. By 21 July 1847, 2283 factorizations had been entered in tables in the order reproduced here. The remaining 169 were calculated later, and I have included them in this printing, whose arrangement differs from the original only in using repeated factors instead of powers to make typesetting easier.
\par The manuscripts containing these last calculations seem to give the results in the form in which they were first found. The order in which the numbers $a$ appear suggests that the divisors of $aa+1$ and so on were sought only for the smaller numbers, and that the larger ones were obtained from these by particular devices. Only the following rule is recorded, however: \textit{From three numbers $a$, $2a-n$, $2a+n$ a fourth is found}
\[
\frac{4a^3-(nn-3)a}{nn+1}
\]
\textit{This is always an integer for $n=0$ and $n=1$, otherwise only}
\[
\begin{aligned}
&\textit{for}\quad a\equiv0\quad\textit{and}\quad\equiv\pm\surd-1\quad mod(nn+1)\quad\textit{when }n\textit{ is even}\\
&\textit{and for}\quad a\equiv0\quad\textit{and}\quad\equiv\pm\surd-1\quad mod\frac{nn+1}{2}\quad\textit{when }n\textit{ is odd}
\end{aligned}
\]
\begin{center}
\begin{tabular}{lrrr}
\textit{Examples} & $a=253,$ & $n=6,$ & $1750507$\\
 & $a=294,$ & $n=11,$ & $832902$\\
 & $a=119,$ & $n=1,$ & $3370437$\\
 & $a=57,$ & $n=3,$ & $74043$\\
 & $a=123,$ & $n=9,$ & $90657$
\end{tabular}
\end{center}
\par The fourth number has no prime divisors other than those of the first three, and of these only the odd prime divisors belonging to the number $n$ are excluded.
\par The table ‘\textit{Quadratorum myrias prima}’ gives the thousands and hundreds in its heading row, the tens and units of the underlying number in its first vertical column, and, in the last
